\subsection{Declarations}

\texttt{<decl-type-specifier> <declarator>}
\begin{mdframed}[nobreak=true,backgroundcolor=gray!10,linecolor=Firebrick4]
\textbf{Examples:}\\
\texttt{int x; // <int> <x>}\\
\texttt{int *x; // <int> <*x>}\\
\texttt{const int *x; // <const int> <*x>}\\
\texttt{int *const x; // <int> <const *x>}\\
\texttt{const int *const x; // <const int> <const *x>}
\end{mdframed}

\subsubsection{Declarator}
Specifies the name being declared and modifies the base type via operators (\texttt{*}, \texttt{\&}, \texttt{[]}, \texttt{()}, etc.). May include cv-qualifiers (\texttt{const}, \texttt{volatile}) in certain positions.

Eg: \texttt{int *const p;} --- \texttt{const} is part of the declarator, yielding "const pointer to int".

\subsubsection{\texttt{decl-type-specifiers}}

\textbf{Fundamental types}\\
\begin{tabularx}{\linewidth}{lX}
    \hline
    \texttt{char}       & \\
    \texttt{int}        & \\
    \texttt{float}      & \\
    \texttt{double}     & \\
    \texttt{void}       & \\
    \hline
Common aliases & \\
    \hline
    \texttt{size\_t}     & \\
    \texttt{char\_t}     & \\
    \texttt{int8\_t}     & \\
    \texttt{int64\_t}    & \\
    \texttt{uint32\_t}   & \\
    \hline
\end{tabularx}
\textcolor{magenta}{NOTE:} \texttt{std::size\_t}, \texttt{int8\_t} are
preferred over \texttt{size\_t}, \texttt{int8\_t} etc. \textcolor{red}{Why?}


\begin{tabularx}{\linewidth}{lX}
    \hline
    \texttt{register}       & Ignored by modern compilers.\\ 
    \texttt{mutable}        & Removes const-ness from a const object member variable.\\
    \texttt{inline}         & Relaxes ODR.\\
                            & \textcolor{red}{NOTE:} Global linkage.\\
    \hline
\end{tabularx}

\begin{tabularx}{\linewidth}{lX}
    \texttt{static}         & Specifies internal linkage.\\
                            & \textit{Local:} initialized only once, retains value between calls.\\
                            & \textit{Global:} visible only in the current TU.\\
                            & \textcolor{red}{Different from \texttt{inline}}.\\
                            & \textit{See "Class" for static members.}\\
    \texttt{extern}         & Specifies external linkage.\\
                            & Eg: \texttt{extern int x; // x is now available in other TUs.}\\
    \texttt{thread\_local}  & Specifies thread local storage duration.\\
                            & Eg: \texttt{thread\_local int x;}\\
    \hline
\end{tabularx}

\vfill\null
\columnbreak

\begin{tabularx}{\linewidth}{lX}
    \texttt{typedef}        & \texttt{typedef <type> <alias>;}\\
                            & Eg: \texttt{typedef long long int ll;}\\
    \texttt{using}          & \texttt{using <alias> = <type>;}\\
                            & Eg: \texttt{using ll = long long int;}\\
                            & \texttt{using} supports template aliase, but \texttt{typedef} does not.\\
    \texttt{namespace}      & Declares a namespace.\\
                            & Eg: \texttt{namespace Foo \{ ... \}}\\
                            & Alias: \texttt{namespace Bar = Foo;}\\
    \hline
    \texttt{auto}          & Type deduction based on initializer.\\
                           & Eg: \texttt{auto x = 10; // x is int}\\
    \texttt{decltype}      & Deduces type of an expression.\\
                           & Eg: \texttt{decltype(x) y; // y has same type as x}\\
                           & \texttt{decltype} is exact, unlike \texttt{auto}.\\
                           & Eg: \texttt{auto} doesn't preserve references, const, arrays decay to pointers etc.\\
    \hline
\end{tabularx}

\begin{tabularx}{\linewidth}{lX}
    \texttt{static\_assert} & Compile time assertion.\\
                            & \texttt{static\_assert(<bool-constexpr>)}\\
                            & \texttt{static\_assert(<bool-constexpr>, "<message>")}\\
    \hline
\end{tabularx}

\subsubsection{enum, enum class, union}

\textbf{enum}
Eg:
\begin{mdframed}[nobreak=true,backgroundcolor=gray!10,linecolor=Firebrick4]
\begin{verbatim}
enum Color { RED,       // 0
             GREEN = 5, // 5
             BLUE       // 6
};
\end{verbatim}
\end{mdframed}

\begin{itemize}
\item No scope:Eg: \texttt{Color c = RED;}.
\item Implicitly converts to int. Eg: \texttt{int x = RED;}.
\end{itemize}

\textbf{enum class (scoped enum)}
Eg:
\begin{mdframed}[nobreak=true,backgroundcolor=gray!10,linecolor=Firebrick4]
\begin{verbatim}
enum class Color {
    Red,        // 0
    Green = 3,  // 3
    Blue        // 4
};
\end{verbatim}
\end{mdframed}

Eg usage: \texttt{Color c = Color::Red;}.\\

\textbf{Union}

Eg:
\begin{mdframed}[nobreak=true,backgroundcolor=gray!10,linecolor=Firebrick4]
\begin{verbatim}
union Data {
    int i;
    float f;
};
\end{verbatim}
\end{mdframed}
